% มคอ.5 — 4095102 ภาคเรียนที่ 1/2568 (รุ่นแรก)
\documentclass{mko5}
\begin{document}
\mkofivetitle

\coursecode{4095102}
\coursename{การเขียนโปรแกรมขั้นสูงสำหรับการเรียนรู้ของเครื่อง}
\curriculumnote{หลักสูตรวิศวกรรมศาสตรมหาบัณฑิต สาขาวิชาวิศวกรรมคอมพิวเตอร์และปัญญาประดิษฐ์ พ.ศ. 2568}
\coursegroup{วิชาเฉพาะด้านบังคับ}
\coursecredits{3 (2-2-5)}
\studyplan{แผน 2 --- วิชาชีพ (สารนิพนธ์)}
\semester{ภาคเรียนที่ 1 / ปีการศึกษา 2568}
\prerequisite{ไม่มี}
\corequisite{ไม่มี}
\studyplace{มหาวิทยาลัยราชภัฏอุตรดิตถ์ --- คณะเทคโนโลยีอุตสาหกรรม}
\registeredstudents{4 คน}
\completedstudents{4 คน}
\instructor{ผศ.ดร.พิศิษฐ์ นาคใจ และ ผศ.ดร.พัชรี มณีรัตน์}
\coinstructors{---}

\mkoheaderinfo

\begin{teachingactualtable}
  \actualrow{CLO1}{บรรยาย Python/ไลบรารี ML สาธิต ทดลองตามตัวอย่าง แบบทดสอบความเข้าใจ}
  \actualrow{CLO2}{Workshop การเตรียมและทำความสะอาดข้อมูล โปรเจกต์กลุ่ม แบบฝึกหัดการจัดการข้อมูล}
  \actualrow{CLO3}{สอนการสร้างแผนภาพ สาธิตการตีความผล รายงานการวิเคราะห์ข้อมูล}
  \actualrow{CLO4}{อภิปรายจริยธรรม กรณีศึกษา รายงานวิเคราะห์จริยธรรม}
  \actualrow{CLO5}{สาธิต Git/GitHub แนะนำ deploy model โครงงาน ML นำเสนอผลงาน}
\end{teachingactualtable}

\begin{assessmentusedtable}
  \assessusedrow{CLO1}{แบบทดสอบความเข้าใจ การนำเสนอและอธิบายหลักการ การตอบคำถาม}
  \assessusedrow{CLO2}{แบบฝึกหัดการจัดการข้อมูล โปรเจกต์กลุ่ม การนำเสนอผลการวิเคราะห์}
  \assessusedrow{CLO3}{รายงานการวิเคราะห์ข้อมูล คุณภาพแผนภาพ การนำเสนอและอภิปรายผล}
  \assessusedrow{CLO4}{รายงานวิเคราะห์จริยธรรม การมีส่วนร่วมในการอภิปราย}
  \assessusedrow{CLO5}{คุณภาพโค้ดและ version control ประสิทธิภาพแบบจำลอง การนำเสนอโครงงาน}
\end{assessmentusedtable}

\mkosectionthree

\closummaryblock{CLO1}{อธิบาย หลักการของภาษาโปรแกรมสมัยใหม่และไลบรารีสำหรับการเรียนรู้ของเครื่อง ได้อย่างถูกต้อง}{80}
\achievementpass{4}{100}{ผ่านเกณฑ์ระดับ 3 ขึ้นไป}
\achievementfail{0}{0}{ไม่มี}
\closummaryend

\closummaryblock{CLO2}{ประยุกต์ใช้ เทคนิคการจัดการข้อมูล ตั้งแต่การเตรียม ทำความสะอาด แปลง และแบ่งชุดข้อมูลสำหรับแบบจำลองการเรียนรู้ของเครื่อง}{80}
\achievementpass{4}{100}{โปรเจกต์กลุ่มและแบบฝึกหัดผ่าน}
\achievementfail{0}{0}{ไม่มี}
\closummaryend

\closummaryblock{CLO3}{วิเคราะห์ ข้อมูลด้วยแผนภาพและการตีความ เพื่อเลือกใช้และปรับปรุงแบบจำลองที่เหมาะสม}{80}
\achievementpass{4}{100}{รายงานวิเคราะห์ข้อมูลผ่านเกณฑ์}
\achievementfail{0}{0}{ไม่มี}
\closummaryend

\closummaryblock{CLO4}{แสดงออก ถึงความตระหนักในจริยธรรมและความรับผิดชอบต่อสังคม ในการพัฒนาโปรแกรมสำหรับการเรียนรู้ของเครื่อง}{80}
\achievementpass{4}{100}{รายงานและการอภิปรายผ่าน}
\achievementfail{0}{0}{ไม่มี}
\closummaryend

\closummaryblock{CLO5}{พัฒนา แบบจำลองการเรียนรู้ของเครื่อง โดยใช้ระบบควบคุมเวอร์ชันและนำไปใช้งานในสภาพแวดล้อมจริงได้อย่างมีประสิทธิภาพ}{80}
\achievementpass{4}{100}{โครงงาน ML และการ deploy ผ่าน}
\achievementfail{0}{0}{ไม่มี}
\closummaryend

\mkosectionfour
\begin{mkotextblock}
การ deploy model ในสภาพแวดล้อม production ต้องการเวลาฝึกเพิ่ม;
ความหลากหลายระดับพื้นฐาน Python ของนักศึกษาไม่เท่ากัน
\end{mkotextblock}

\mkosectionfive
\begin{mkotextblock}
จัดทำ Repository ตัวอย่างพร้อม CI สำหรับโครงงาน;
เพิ่ม Workshop การใช้ MLOps เบื้องต้นในภาคเรียนถัดไป
\end{mkotextblock}

\mkosectionsix
\begin{mkotextblock}
ขอสนับสนุน Cloud credits หรือเซิร์ฟเวอร์สำหรับการ deploy โมเดลของนักศึกษา
\end{mkotextblock}

\mkosectionseven
\mkosignaturepair{ผศ.ดร.พิศิษฐ์ นาคใจ และ ผศ.ดร.พัชรี มณีรัตน์}{30 เมษายน 2569}{ผศ.ดร.พิทักษ์ คล้ายชม}{30 เมษายน 2569}
\end{document}
